%==============================================================
%  Figura 7.1 - Attivita' di verifica e validazione (V&V) della
%  EN ISO 13849-1: estratto del processo iterativo di progettazione
%  Adattamento in italiano della Figura 7.1 dell'IFA Report 2/2017
%  (stesso stile delle Figure 5.5 e 6.1, di cui e' la prosecuzione)
%==============================================================
\begin{tikzpicture}[
  font=\normalsize,
  grigio/.style={fill=black!10, align=center, inner sep=4pt},
  numcella/.style={draw=black, fill=white, line width=0.8pt,
                   minimum width=0.7cm, minimum height=0.75cm, inner sep=1pt},
  corpo/.style={draw=black, fill=black!10, line width=0.8pt, align=center,
                minimum width=8.5cm, text width=8.2cm, inner sep=4pt},
  decis/.style={diamond, draw=black, fill=black!10, line width=0.8pt, align=center,
                inner sep=1pt, aspect=3.4, minimum width=8.5cm},
  fdline/.style={draw=black, line width=0.9pt},
  fdar/.style={fdline, -{Triangle[length=6pt,width=6pt]}}
]

%--------------------------------------------------------------
% Fase 5: valutazione del PL (due righe)
%--------------------------------------------------------------
\node[corpo, minimum height=1.4cm, anchor=north] (fase5a) at (7.75,-2.40)
  {Valutazione del PL degli SRP/CS in termini di categoria,
   \textit{MTTF}\textsubscript{D}, \textit{DC}\textsubscript{avg}, CCF};
\node[corpo, minimum height=0.8cm, anchor=north] (fase5b) at (fase5a.south)
  {Software e guasti sistematici};
\node[numcella, inner sep=0pt,
      fit={([xshift=-0.7cm]fase5a.north west) (fase5b.south west)}] (num5) {5};

%--------------------------------------------------------------
% Fasi 6, 7 e 8: verifica, validazione e completezza dell'analisi
%--------------------------------------------------------------
\node[decis, anchor=north] (dec6) at ([yshift=-0.45cm]fase5b.south)
  {Verifica:\\[2pt] PL $\geq$ PL\textsubscript{r}?};
\node[decis, anchor=north] (dec7) at ([yshift=-0.80cm]dec6.south)
  {Validazione:\\[2pt] requisiti soddisfatti?};
\node[decis, anchor=north] (dec8) at ([yshift=-0.80cm]dec7.south)
  {Tutte le SF\\[2pt] analizzate?};

\node[numcella, anchor=east] (num6) at ([xshift=-0.12cm]dec6.west) {6};
\node[numcella, anchor=east] (num7) at ([xshift=-0.12cm]dec7.west) {7};
\node[numcella, anchor=east] (num8) at ([xshift=-0.12cm]dec8.west) {8};

%--------------------------------------------------------------
% Ingresso dalla realizzazione e determinazione del PL
%--------------------------------------------------------------
\draw[fdar] (7.75,-1.85) -- (fase5a.north);
\node[align=center, anchor=south] (ingresso) at (7.75,-1.80)
  {Dall'implementazione\\ e determinazione del PL\\ (Figura 6.1)};

%--------------------------------------------------------------
% Percorso principale
%--------------------------------------------------------------
\draw[fdar] (fase5b.south) -- (dec6.north);
\draw[fdar] (dec6.south) -- (dec7.north);
\draw[fdar] (dec7.south) -- (dec8.north);
\node[anchor=east] at ([xshift=-0.08cm,yshift=0.40cm]dec7.north) {s\`i};
\node[anchor=east] at ([xshift=-0.08cm,yshift=0.40cm]dec8.north) {s\`i};

%--------------------------------------------------------------
% "Per ciascuna SF" con graffa
%--------------------------------------------------------------
\coordinate (bx) at ([xshift=-0.35cm]num5.west);
\coordinate (bgiu) at (bx |- dec7.south);
\coordinate (balto) at (bx |- fase5a.north);
\coordinate (bmezzo) at (bx |- dec6.south);
\draw[fdline, decorate, decoration={brace, amplitude=6pt}] (bgiu) -- (balto);
\node[grigio, minimum width=1.5cm, minimum height=1.9cm, anchor=east]
  (perogni) at ([xshift=-0.25cm]bmezzo) {Per\\ ciascuna\\ SF};

%--------------------------------------------------------------
% Rami "no": ritorno alla riprogettazione e iterazione sulle SF
%--------------------------------------------------------------
\coordinate (r6) at (13.20,0 |- dec6.east);
\draw[fdline] (dec6.east) -- (r6);
\draw[fdar]   (r6) -- (13.20,-1.85);
\node[anchor=south east] at ([xshift=-0.05cm]r6) {no};

\coordinate (r7) at (13.55,0 |- dec7.east);
\draw[fdline] (dec7.east) -- (r7);
\draw[fdar]   (r7) -- (13.55,-1.85);
\node[anchor=south east] at ([xshift=-0.05cm]r7) {no};

\coordinate (r8) at (15.50,0 |- dec8.east);
\draw[fdline] (dec8.east) -- (r8);
\draw[fdar]   (r8) -- (15.50,-1.85);
\node[anchor=south east] at ([xshift=-0.05cm]r8) {no};

%--------------------------------------------------------------
% Blocco "Verifica" (sovrapposto ai rami di ritorno)
%--------------------------------------------------------------
\node[draw=black, fill=white, line width=0.8pt, align=center,
      minimum width=3.6cm, minimum height=1.3cm] (verif) at (14.35,-2.97) {Verifica};

\node[align=center, anchor=south] at (12.35,-1.80)
  {Alla riprogettazione\\ e nuova\\ determinazione del PL\\ (Figura 6.1)};
\node[align=center, anchor=south] at (15.60,-1.80)
  {Iterazione:\\ altre SF\\ (Figura 5.5)};

%--------------------------------------------------------------
% Uscita verso l'analisi dei rischi
%--------------------------------------------------------------
\coordinate (giu) at ([yshift=-0.60cm]dec8.south);
\node[grigio, minimum width=4.2cm, minimum height=1.3cm]
  (analisi) at (2.35,0 |- giu) {All'analisi dei rischi\\ (EN ISO 12100)};
\draw[fdline] (dec8.south) -- (giu);
\draw[fdar]   (giu) -- (analisi.east);
\node[anchor=east] at ([xshift=-0.08cm,yshift=-0.30cm]dec8.south) {s\`i};

%--------------------------------------------------------------
% Cornice
%--------------------------------------------------------------
\draw[draw=black, line width=1pt]
  ([shift={(-0.45,0.45)}]current bounding box.north west) rectangle
  ([shift={(0.45,-0.45)}]current bounding box.south east);

\end{tikzpicture}
